Providers of large language models (LLMs) retire model versions on fixed dates, and applications that still call a retired model stop working. A recent study by \citet{kim2026retires} reported that \KimAfterShutdown\% of open-source migrations away from a retired model happen only after its shutdown. We replicate and extend that study on a new dataset of \MigrationsAway{} commits from public GitHub repositories that move away from one of \ModelIds{} retired OpenAI, Anthropic and Google model IDs, with a migration label whose precision we measured on a blind, stratified sample (\Precision\%, 95\% CI \PrecisionLow--\PrecisionHigh). For models with exact retirement dates and at least 90 days of follow-up, \AfterLate\% (95\% CI \AfterLateLow--\AfterLateHigh) of migrations come after the shutdown: the direction holds, but the share is clearly below \KimAfterShutdown\%, and Kim's own labelling rules give \AfterKimLabel\% on our data. The median migration lands \ExactDays{} days after the shutdown, the busiest window is the first month after it, and projects that miss the deadline wait longer for OpenAI models than for Anthropic ones. \FollowsLine\% of migrations with a named target follow the provider's recommended model line, yet \TargetOnList\% move to a model that is itself already scheduled for retirement. Only \TouchTest\% of sampled migration commits touch a test file, in line with Kim's \KimTouchTest\%, and \TouchEval\% touch an evaluation file. The results hold when commits are grouped by repository, when other labels are used, and when search gaps are weighted back in. We release the dataset and the code.
